\section{Experimental Methodology and Proposed Evaluation Protocol}
\label{sec:eval}
\begin{table*}[!t]
\centering
\caption{Proposed Evaluation Protocol. None of These Experiments Has Been
Run; They Are Specified Here as the Required Next Step, Not Reported as
Results.}
\label{tab:evalprotocol}
\footnotesize
\renewcommand{\arraystretch}{1.25}
\begin{tabular}{@{}p{2.5cm}p{3.4cm}p{4.0cm}p{4.7cm}@{}}
\toprule
\textbf{Experiment} & \textbf{Metrics} & \textbf{Data / apparatus} & \textbf{Baselines / conditions} \\
\midrule
E1 Retrieval quality &
P@1, P@3, P@5, R@5, MRR, nDCG@10 \cite{jarvelin2002ndcg} &
$\ge$150 institution queries with graded relevance judgments over the
$1{,}507$-chunk corpus &
BM25 only; dense only; hybrid $\lambda\!\in\!\{0.3,0.5,0.6,0.7\}$; effect of
query expansion \\

E2 Reranker ablation &
$\Delta$ nDCG@5, P@5 vs.\ no rerank &
E1 query set + feature vectors &
no rerank; heuristic (Eq.~\eqref{eq:rerank}); trained SVR (after labeling);
weight sensitivity \\

E3 Intent classification &
accuracy, macro/weighted $F_1$, per-class $F_1$, $12{\times}12$ confusion &
$\ge$600 real logged queries, human-labeled, stratified split &
current LSTM (as-is); LSTM retrained on real data; keyword-rule router;
zero-shot LLM classifier \\

E4 Answer faithfulness &
faithfulness, answer relevancy, context relevancy \cite{es2024ragas};
manual hallucination rate; citation correctness &
$\ge$120 questions (in-domain / out-of-scope / ambiguous); LLM judge +
2 human raters, report $\kappa$ &
full pipeline; ablate confidence gate; ablate grounding check; ablate both \\

E5 Confidence calibration &
coverage, refusal rate, hallucination rate, answered-accuracy vs.\
$(\theta_{\text{MED}}, \theta_{\text{HIGH}})$; reliability diagram, ECE &
E4 question set with correctness labels &
sweep $\theta$ on a grid; with vs.\ without LSTM $P(\text{intent})$ in
Eq.~\eqref{eq:conf} \\

E6 Pathfinding &
path length, wall-clock, expanded-node count vs.\ graph size &
the $180$/$325$ graph + synthetically scaled graphs; synthetic coordinates
for tier-1/2 heuristic &
A$^{\star}$ (3 heuristic tiers) vs.\ Dijkstra; cached vs.\ rebuilt graph \\

E7 Tour traversal &
cross-floor transition success, guided-tour completion rate, orientation
continuity error, interaction errors &
scripted traversals over all $6$ floor buckets; fault-injected corrupt
graphs &
manual vs.\ guided; with vs.\ without cycle guard \\

E8 Usability &
task completion rate, time-on-task, SUS, satisfaction, trust &
$\ge$20 participants (prospective students), $6$ tasks spanning QA + tour +
(future) routing &
\system{} vs.\ institution website + PDF brochures \\
\bottomrule
\end{tabular}
\end{table*}


The current artifact contains \emph{no} formal evaluation pipeline: there
are no retrieval precision/recall figures, no answer-faithfulness scores,
no reranker comparison, no full intent-classifier confusion matrix beyond
the one reproduced in Section~\ref{sec:results}, no load-test results, no
pathfinding benchmark, and no user study. Rather than fabricate any of
these, this section specifies exactly how each should be measured. The
protocol is summarized in Table~\ref{tab:evalprotocol}; the subsections
below give the design detail. We distinguish \emph{implementation metrics}
(Table~\ref{tab:counts}, already measured) from \emph{experimental
performance metrics} (E1--E8 below, not yet measured).

\subsection{E1 --- Retrieval Quality}
\emph{Objective:} test whether hybrid fusion~\eqref{eq:fusion} outperforms
its components on this corpus, and whether $\lambda = 0.6$ is well chosen.
\emph{Apparatus:} construct a query set of at least $150$ realistic
institution questions (admissions, academics, facilities, campus info,
out-of-scope) and, for each, human-annotate graded relevance for the pooled
top-$k$ of BM25, dense, and hybrid retrieval, following BEIR-style pooling
practice~\cite{thakur2021beir}. \emph{Metrics:} precision@$\{1,3,5\}$,
recall@$5$, MRR, and nDCG@$10$~\cite{jarvelin2002ndcg}. \emph{Conditions:}
BM25 only; dense only; hybrid at
$\lambda \in \{0.3, 0.5, 0.6, 0.7\}$; hybrid with and without query
expansion. \emph{Expected finding form:} ``hybrid at $\lambda=x$ achieves
nDCG@10 $= y$, a $\Delta$ of $z$ over the stronger single method
($p < 0.05$, paired bootstrap).'' No such claim is made here.

\subsection{E2 --- Reranker Ablation}
\emph{Objective:} quantify the heuristic reranker's contribution and, once
E1 has produced relevance labels, train and evaluate the currently-untrained
SVR. \emph{Conditions:} no reranking; the four-feature heuristic of
Eq.~\eqref{eq:rerank}; $\varepsilon$-SVR trained on E1 labels; a sensitivity
sweep of the four heuristic weights. \emph{Metrics:} $\Delta$nDCG@$5$ and
$\Delta$P@$5$ against the no-rerank baseline. Until the SVR is trained, no
number for it may be reported.

\subsection{E3 --- Intent Classification}
\emph{Objective:} measure generalization on \emph{real} queries, not the
synthetic seed set. \emph{Apparatus:} sample at least $600$ queries from
the accumulated chat logs, have two annotators label them into the $12$
classes (report inter-annotator $\kappa$), and split stratified by class.
\emph{Conditions:} the current LSTM as-is; the LSTM retrained on the real
labeled data; the deterministic keyword router alone; a zero-shot LLM
classifier. \emph{Metrics:} accuracy, macro and weighted $F_1$, per-class
$F_1$, and the full $12{\times}12$ confusion matrix. \emph{Framing:} the
current $0.529$ validation accuracy on $34$ synthetic examples
(Section~\ref{sec:results}) is explicitly \emph{not} a generalization
estimate; E3 is what would produce one.

\subsection{E4 --- Answer Faithfulness and Relevancy}
\emph{Objective:} measure whether generated answers are entailed by their
cited context and whether the safeguards matter. \emph{Apparatus:} a fixed
question set of at least $120$ items spanning in-domain, out-of-scope, and
deliberately ambiguous queries; scoring by an LLM judge for faithfulness,
answer relevancy, and context relevancy~\cite{es2024ragas}, cross-checked
by two human raters on a subset (report agreement); a manually adjudicated
hallucination rate; and citation-correctness (does each cited source
actually support the answer). \emph{Ablations:} full pipeline; remove the
confidence gate; remove the post-generation grounding check; remove both.
\emph{Correct-refusal behavior:} for out-of-scope and ambiguous items,
score whether the system refused or asked a clarifying question rather than
answering. This is the single highest-value missing measurement for a
results section.

\subsection{E5 --- Confidence-Threshold Calibration}
\emph{Objective:} replace the nine-query threshold derivation
(Table~\ref{tab:confsample}) with a statistically sized calibration.
\emph{Apparatus:} the E4 question set with per-item correctness labels.
\emph{Analysis:} sweep $(\theta_{\text{MED}}, \theta_{\text{HIGH}})$ on a
grid and plot coverage, refusal rate, hallucination rate, and
answered-accuracy; produce a reliability diagram and report expected
calibration error. \emph{Additional condition:} repeat with the LSTM's
$P(\text{intent})$ actually supplied to Eq.~\eqref{eq:conf} (closing the
implementation gap noted in Section~\ref{sec:rag}), to test whether the
intent term improves separation.

\subsection{E6 --- Pathfinding Performance}
\emph{Objective:} characterize A$^{\star}$ on the live graph and against
Dijkstra. \emph{Apparatus:} the $180$-node/$325$-edge graph plus
synthetically scaled graphs (up to $10^4$ nodes) and synthetic node
coordinates so tiers~1 and~2 of the heuristic~\eqref{eq:heur} are
exercised. \emph{Metrics:} route length (should be identical to Dijkstra,
verifying optimality), wall-clock time, and expanded-node count as a
function of graph size and heuristic tier; the marginal cost of rebuilding
the graph per request vs.\ caching it.

\subsection{E7 --- Tour Traversal Correctness}
\emph{Objective:} verify the scene-graph mechanisms. \emph{Apparatus:}
scripted traversals covering every one of the six floor buckets and every
cross-floor boundary; fault-injected corrupt graphs (self-loops, dangling
tails). \emph{Metrics:} cross-floor transition success rate, guided-tour
completion rate, camera-orientation continuity error at transitions
(deviation from the Street-View heading-preservation rule), and interaction
errors; verification that the cycle guard aborts to \emph{error} rather
than looping.

\subsection{E8 --- Usability}
\emph{Objective:} measure whether the integrated system helps real users.
\emph{Apparatus:} at least $20$ prospective-student participants performing
six tasks that each require both propositional and spatial information
(e.g.\ ``find which floor the Power System Simulation Lab is on and view
it''). \emph{Comparison:} \system{} vs.\ the institution's website plus PDF
brochures. \emph{Metrics:} task completion rate, time-on-task, System
Usability Scale, satisfaction, and a trust item (``I would rely on this
system's answers''). \emph{Analysis:} paired tests with effect sizes.

\subsection{Reproducibility Requirements}
Any reported result must also record: exact hardware (CPU/GPU, RAM, whether
Ollama ran on CPU or GPU); the Llama~3.2 quantization and Ollama version;
the corpus snapshot (chunk count and manifest hash); and the random seeds.
The current paper reports none of these because no experiment has been run.
